% ==========================================================
% fariscovernew.sty — دليل الاستخدام الكامل
% Compile: xelatex
% ==========================================================
\documentclass[a4paper,12pt]{report}
\usepackage[margin=1.5cm]{geometry}

% 1-part-Languages and Fonts
\usepackage{fontspec}
\usepackage{polyglossia}
\setmainlanguage[numerals=eastern]{arabic}
\setotherlanguage{english}
\newfontfamily\arabicfont[Script=Arabic]{Amiri}
\newfontfamily\englishfont{Latin Modern Roman}

% 2-part-Mono font
\setmonofont{DejaVu Sans Mono}[Scale=MatchLowercase]

% 3-part-amsmath
\usepackage{amsmath}

\usepackage{xcolor}
\usepackage{listings}

% 4-part-Safe Arabic macro inside code
\newcommand{\AR}[1]{\(\text{\beginR\arabicfont #1\endR}\)}
\newcommand{\en}[1]{\textenglish{ #1}}
% 5-part-Listing style
\lstdefinestyle{fariscode}{
  basicstyle=\ttfamily\small,
  breaklines=true,
  frame=single,
  rulecolor=\color{gray!60},
  backgroundcolor=\color{gray!3},
  columns=fullflexible,
  keepspaces=true,
  language={[LaTeX]TeX},
  showstringspaces=false,
  escapeinside={@}{@}
}

% 6-part-Package
\usepackage{fariscovernew}

% 7-part-Helpers
\usepackage{adjustbox}
\usepackage{fancybox}

\title{دليل استعمال الحزمة (\textenglish{fariscovernew - v1.1}) باللغة العربية\\
مع إمكانية عرض الأوامر المختلفة والكود}
\author{تأليف الأستاذ فارس الحيدري}
\date{\en{2026/03/21}}

\makeatletter
\def\@seccntformat#1{\csname the#1\endcsname\quad}
\makeatother

\usepackage[unicode,colorlinks=true,linkcolor=blue,bookmarks=false]{hyperref}

\pdfstringdefDisableCommands{%
  \def\textarabic#1{#1}%
  \def\textenglish#1{#1}%
}

\usepackage{chngcntr}

\AtBeginDocument{%
  \counterwithout{section}{chapter}%
  \counterwithin{subsection}{section}%
  \counterwithin{subsubsection}{subsection}%
  \setcounter{secnumdepth}{3}%
  \setcounter{tocdepth}{3}%
  \makeatletter
  \renewcommand\thesection{\arabic{section}}%
  \renewcommand\thesubsection{\arabic{section}-\arabic{subsection}}%
  \renewcommand\thesubsubsection{\arabic{section}-\arabic{subsection}-\arabic{subsubsection}}%
  \makeatother
}

% ====================================

\begin{document}

\maketitle
\newpage
\tableofcontents

\clearpage

% ============================ section= one ==========================
\section{سطر واحد بسيط لتصميم صفحة عنوان مميزة}
% ==================================================================

\begin{english}
\begin{lstlisting}[style=fariscode]
\FarisTitlePage[A][90][45][RoyalBlue][Amiri][Amiri]{ @\AR{العنوان الرئيسي}@ }{ @\AR{العنوان الفرعي}@ }
\end{lstlisting}
\end{english}

\vspace{0.5cm}

\noindent\textbf{\AR{النتيجة}:}\\[0.3cm]
\adjustbox{scale=0.3,center}{%
\begin{minipage}{\textwidth}
\FarisTitlePage[A][90][60][RoyalBlue][Amiri][Amiri]{العنوان الرئيسي}{العنوان الفرعي}
\end{minipage}
}

\vspace{6cm}

% ============================ section= second ========================
\section{الأنماط الأساسية التي تحتويها الحزمة}
% ==================================================================

\begin{english}
\begin{lstlisting}[style=fariscode]
\FarisSetStyle{minimal}   % @\AR{النمط البسيط}@
\FarisSetStyle{triple}    % @\AR{نمط ثلاثي}@
\FarisSetStyle{diamond}   % @\AR{النمط الماسي البسيط}@
\FarisSetStyle{diamondx}  % @\AR{النمط الماسي المتقدم}@
\FarisSetStyle{wave}      % @\AR{نمط متموج}@
\FarisSetStyle{sidebar}   % @\AR{شريط جانبي}@
\FarisSetStyle{diagonal}  % @\AR{النمط القطري}@
\FarisSetStyle{badge}     % @\AR{نمط الشارة}@
\FarisSetStyle{ribbon}    % @\AR{النمط الشريطي}@
\end{lstlisting}
\end{english}

\subsection{أوامر تخصيص خاصة بـ \textenglish{diamondx}}

\begin{english}
\begin{lstlisting}[style=fariscode]
% @\AR{الألوان المتممة}@
\FarisDiamondBackComplementOn
\FarisDiamondBackComplementOff

% @\AR{التينتات}@
\FarisSetDiamondTints{front}{backA}
\FarisSetDiamondBackTintII{backB}

% @\AR{الزوايا}@
\FarisSetDiamondAnglePair{angle}
\FarisSetDiamondAngles{backAngle}{frontAngle}
\FarisSetDiamondBackDelta{delta}

% @\AR{الألوان الثلاثية}@
\FarisSetDiamondTriadOffset{offset}
\end{lstlisting}
\end{english}

\newpage

% ============================ section= Three =========================
\section{أمثلة سريعة لكل نمط}
% ==================================================================

\begin{english}
\begin{lstlisting}[style=fariscode]
% @\AR{ثلاثي}@
\FarisSetStyle{triple}
\FarisTitlePage[A][70][-][Blue][Amiri][Amiri]{@\AR{العنوان}@}{@\AR{الفرعي}@}

% @\AR{قطري}@
\FarisSetStyle{diagonal}
\FarisTitlePage[A][70][-][Red][Amiri][Amiri]{@\AR{العنوان}@}{@\AR{الفرعي}@}

% @\AR{موجي}@
\FarisSetStyle{wave}
\FarisTitlePage[A][70][-][Green][Amiri][Amiri]{@\AR{العنوان}@}{@\AR{الفرعي}@}

% @\AR{شارة}@
\FarisSetStyle{badge}
\FarisTitlePage[A][70][-][Purple][Amiri][Amiri]{@\AR{العنوان}@}{@\AR{الفرعي}@}

% @\AR{ماسي بسيط}@
\FarisSetStyle{diamond}
\FarisTitlePage[A][70][-][Orange][Amiri][Amiri]{@\AR{العنوان}@}{@\AR{الفرعي}@}

% @\AR{بسيط}@
\FarisSetStyle{minimal}
\FarisTitlePage[A][70][-][Gray][Amiri][Amiri]{@\AR{العنوان}@}{@\AR{الفرعي}@}

% @\AR{شريطي}@
\FarisSetStyle{ribbon}
\FarisTitlePage[A][70][-][Maroon][Amiri][Amiri]{@\AR{العنوان}@}{@\AR{الفرعي}@}

% @\AR{جانبي}@
\FarisSetStyle{sidebar}
\FarisTitlePage[A][70][-][Navy][Amiri][Amiri]{@\AR{العنوان}@}{@\AR{الفرعي}@}

% @\AR{ماسي متقدم}@
\FarisSetStyle{diamondx}
\FarisSetColor{SlateBlue}
\FarisDiamondBackComplementOn
\FarisSetDiamondTints{30}{40}
\FarisTitlePage[A][70][-][SlateBlue][Amiri][Amiri]{@\AR{العنوان}@}{@\AR{الفرعي}@}
\end{lstlisting}
\end{english}

\newpage

% ============================ section= Four ==========================
\section{مثال كامل لاستخدام \textenglish{diamondx}}
% ==================================================================

\begin{english}
\begin{lstlisting}[style=fariscode]
\FarisSetStyle{diamondx}
\FarisSetColor{SlateBlue}
\FarisDiamondBackComplementOn
\FarisSetDiamondTints{30}{30}      % @\AR{تينت الأمامي و(A)}@
\FarisSetDiamondBackTintII{40}     % @\AR{تينت (B)}@
\FarisSetDiamondAnglePair{45}      % @\AR{زوايا ±45°}@
\FarisSetDiamondBackDelta{15}      % @\AR{B = A + 15°}@
\FarisSetDiamondTriadOffset{75}    % @\AR{لون B بين الأساس والمتمّم}@
\FarisTitlePage[A][90][-][SlateBlue][Amiri][Amiri][0.8]{@\AR{عنوان متقدم}@}{@\AR{عنوان فرعي}@}
\end{lstlisting}
\end{english}

\vspace{2cm}

\noindent\textbf{\AR{النتيجة}:}\\[0.3cm]
\adjustbox{scale=0.3,center}{%
\begin{minipage}{\textwidth}
\FarisSetStyle{diamondx}
\FarisSetColor{SlateBlue}
\FarisDiamondBackComplementOn
\FarisSetDiamondTints{30}{30}
\FarisSetDiamondBackTintII{40}
\FarisSetDiamondAnglePair{45}
\FarisSetDiamondBackDelta{15}
\FarisSetDiamondTriadOffset{75}
\FarisTitlePage[A][90][-][SlateBlue][Amiri][Amiri][0.8]{عنوان متقدم}{عنوان فرعي}
\end{minipage}
}

\newpage

% ============================ section = Five ==========================
\section{تحديد حجم النصوص}
% ==================================================================

\subsection{تحديد حجم النص الأساسي \textenglish{90pt} فقط}

\begin{english}
\begin{lstlisting}[style=fariscode]
% @\AR{العنوان الفرعي يأخذ الحجم التلقائي 4/5 من الرئيسي}@
\FarisTitlePage[A][90][-][DarkOrange1][Diwani Simple Striped][Amiri]{@\AR{الفصل الأول}@}{@\AR{العنوان الفرعي}@}
\end{lstlisting}
\end{english}

\vspace{2cm}

\noindent\textbf{\AR{النتيجة}:}\\[0.3cm]
\adjustbox{scale=0.3,center}{%
\begin{minipage}{\textwidth}
\FarisTitlePage[A][90][-][DarkOrange1][Diwani Simple Striped][Amiri]{الفصل الأول}{العنوان الفرعي}
\end{minipage}
}

\vspace{6cm}

\subsection{تحديد الرئيسي \textenglish{90pt} والفرعي \textenglish{45pt}}

\begin{english}
\begin{lstlisting}[style=fariscode]
% @\AR{تحديد حجم النص الأساسي والعنوان الفرعي}@
% @\AR{ملاحظة: إذا كان الفرعي أكبر من الرئيسي يُعاد تلقائياً لـ 4/5 من الرئيسي}@
\FarisTitlePage[A][90][45][DarkOrange1][Diwani Simple Striped][Amiri]{@\AR{الفصل الأول}@}{@\AR{العنوان الفرعي}@}
\end{lstlisting}
\end{english}

\vspace{6.5cm}

\noindent\textbf{\AR{النتيجة}:}\\[0.3cm]
\adjustbox{scale=0.3,center}{%
\begin{minipage}{\textwidth}
\FarisTitlePage[A][90][45][DarkOrange1][Diwani Simple Striped][Amiri]{الفصل الأول}{العنوان الفرعي}
\end{minipage}
}

\newpage

\subsection{تحديد الرئيسي \textenglish{90pt} والفرعي بالنسبة \textenglish{0.5}}

\begin{english}
\begin{lstlisting}[style=fariscode]
% @\AR{حجم الفرعي = 0.5 × حجم الرئيسي = 45pt}@
\FarisTitlePage[A][90][-][Crimson][Amiri][Amiri][0.5]{@\AR{الفصل الثاني}@}{@\AR{العنوان الفرعي}@}
\end{lstlisting}
\end{english}

\vspace{1cm}

\noindent\textbf{\AR{النتيجة}:}\\[0.3cm]
\adjustbox{scale=0.3,center}{%
\begin{minipage}{\textwidth}
\FarisTitlePage[A][90][-][Crimson][Diwani Simple Striped][Amiri][0.5]{الفصل الثاني}{العنوان الفرعي}
\end{minipage}
}

\vspace{8cm}

\subsection{قواعد حجم النص — جدول مرجعي}

\begin{center}
\begin{tabular}{|c|c|c|c|}
\hline
\textbf{TitlePt} & \textbf{SubtitlePt} & \textbf{SubRatio} & \AR{النتيجة} \\
\hline
90 & - & - & \AR{رئيسي 90 · فرعي
\en{ 72 (4/5)}} \\
\hline
90 & 45 & - & \AR{رئيسي 90 · فرعي 45} \\
\hline
90 & - & 0.5 & \AR{رئيسي 90 · فرعي 45} \\
\hline
0 & 45 & - & \AR{رئيسي \en{56.25 (5/4×45)} · فرعي 45} \\
\hline
- & 45 & - & \AR{رئيسي
\en{56.25 (5/4×45)} · فرعي 45} \\
\hline
90 & 0 & 0.5 & \AR{رئيسي 90 · فرعي 45 (
\en{SubRatio} يأخذ الأولوية)} \\
\hline
90 & 90 & - & \AR{رئيسي 90 · فرعي 72 (حارس التراتبية)} \\
\hline
\end{tabular}
\end{center}

\newpage

% ============================ section = Sex ==========================
\section{تلوين النصوص — الجديد في الإصدار \textenglish{v1.1}}
% ==================================================================

تتيح الحزمة تخصيص لون العنوان الرئيسي والفرعي بشكل مستقل عن لون التصميم العام.
الألوان اختيارية، إذا لم تُحدَّد تُشتق تلقائياً من اللون الأساسي.

\subsection{الأوامر المتاحة}

\begin{english}
\begin{lstlisting}[style=fariscode]
% @\AR{تخصيص لون العنوان الرئيسي}@
\FarisSetTitleTextColor{color}

% @\AR{تخصيص لون العنوان الفرعي}@
\FarisSetSubtitleTextColor{color}

% @\AR{إعادة ضبط لون العنوان الرئيسي للافتراضي}@
\FarisResetTitleTextColor

% @\AR{إعادة ضبط لون العنوان الفرعي للافتراضي}@
\FarisResetSubtitleTextColor

% @\AR{إعادة ضبط كليهما معاً}@
\FarisResetTextColors
\end{lstlisting}
\end{english}

\vspace{0.5cm}

\subsection{قاعدة الأولوية}

\begin{english}
\begin{lstlisting}[style=fariscode]
% @\AR{إذا حُدِّد اللون يدوياً  يأخذ الأولوية دائماً}@
\FarisSetTitleTextColor{orange}    % @\AR{ يُطبَّق}@

% @\AR{إذا كان فارغاً  يُشتق من BaseColor!85!black}@
\FarisResetTitleTextColor         % @\AR{ يعود للافتراضي}@
\end{lstlisting}
\end{english}

\newpage 

\subsection{مثال: عنوان برتقالي على خلفية داكنة}

\begin{english}
\begin{lstlisting}[style=fariscode]
\FarisSetTitleTextColor{orange}
\FarisSetSubtitleTextColor{orange}
\FarisTitlePage[A][90][-][DarkOrange1][Diwani Simple Striped][Amiri]{@\AR{الفصل السادس}@}{@\AR{الكيمياء التحليلية}@}
\end{lstlisting}
\end{english}

\vspace{1cm}

\noindent\textbf{\AR{النتيجة}:}\\[0.3cm]
\adjustbox{scale=0.3,center}{%
\begin{minipage}{\textwidth}
\FarisSetTitleTextColor{orange}
\FarisSetSubtitleTextColor{orange}
\FarisTitlePage[A][90][-][DarkOrange1][Diwani Simple Striped][Amiri]{الفصل السادس}{الكيمياء التحليلية}
\end{minipage}
}

\vspace{6cm}

\subsection{مثال: لون مختلف لكل عنوان}

\begin{english}
\begin{lstlisting}[style=fariscode]
\FarisSetTitleTextColor{NavyBlue}
\FarisSetSubtitleTextColor{DarkOrange1}
\FarisTitlePage[A][90][-][ForestGreen][Amiri][Amiri]{@\AR{الفصل السابع}@}{@\AR{الكيمياء العضوية}@}
\end{lstlisting}
\end{english}

\vspace{6cm}

\noindent\textbf{\AR{النتيجة}:}\\[0.3cm]
\adjustbox{scale=0.3,center}{%
\begin{minipage}{\textwidth}
\FarisSetTitleTextColor{NavyBlue}
\FarisSetSubtitleTextColor{DarkOrange1}
\FarisTitlePage[A][90][-][ForestGreen][Amiri][Amiri]{الفصل السابع}{الكيمياء العضوية}
\end{minipage}
}

\newpage

\subsection{مثال: إعادة الضبط بعد التخصيص}

\begin{english}
\begin{lstlisting}[style=fariscode]
% @\AR{غلاف 1: ألوان مخصصة}@
\FarisSetTitleTextColor{white}
\FarisSetSubtitleTextColor{black}
\FarisTitlePage[A][90][-][DarkOrange1][Amiri][Amiri]{@\AR{الفصل الأول}@}{@\AR{المقدمة}@}

% @\AR{غلاف 2: العودة للافتراضي (مشتق من BaseColor)}@
\FarisResetTextColors
\FarisTitlePage[A][90][-][NavyBlue][Amiri][Amiri]{@\AR{الفصل الثاني}@}{@\AR{الفيزياء}@}
\end{lstlisting}
\end{english}

\vspace{0.5cm}

\subsection{جدول الألوان المدعومة}

\begin{english}
\begin{lstlisting}[style=fariscode]
% @\AR{أسماء xcolor الجاهزة}@
\FarisSetTitleTextColor{white}
\FarisSetTitleTextColor{black}
\FarisSetTitleTextColor{NavyBlue}
\FarisSetTitleTextColor{DarkOrange1}
\FarisSetTitleTextColor{ForestGreen}
\FarisSetTitleTextColor{Crimson}
\FarisSetTitleTextColor{RoyalBlue}
\FarisSetTitleTextColor{Gold1}

% @\AR{ألوان مخلوطة}@
\FarisSetTitleTextColor{blue!70!black}
\FarisSetTitleTextColor{red!80!white}
\FarisSetTitleTextColor{DarkOrange1!60!black}

% @\AR{ألوان RGB مخصصة}@
\definecolor{MyColor}{RGB}{25,118,210}
\FarisSetTitleTextColor{MyColor}
\end{lstlisting}
\end{english}

\newpage

% ============================ section = Seven ========================
\section{الأنماط المتوفرة في الحزمة}
% ==================================================================

\subsection{نمط \textenglish{triple} (النمط الافتراضي)}

\begin{english}
\begin{lstlisting}[style=fariscode]
\FarisSetStyle{triple}
\FarisSetColor{ForestGreen}
\FarisTitlePage[A][80][-][ForestGreen][Amiri][Amiri][0.6]{@\AR{نمط ثلاثي الطبقات}@}{@\AR{النمط الافتراضي للحزمة}@}
\end{lstlisting}
\end{english}

\vspace{6cm}

\noindent\textbf{\AR{النتيجة}:}\\[0.3cm]
\adjustbox{scale=0.3,center}{%
\begin{minipage}{\textwidth}
\FarisResetTextColors
\FarisSetStyle{triple}
\FarisSetColor{ForestGreen}
\FarisTitlePage[A][75][-][ForestGreen][Amiri][Amiri][0.6]{نمط ثلاثي الطبقات}{النمط الافتراضي للحزمة}
\end{minipage}
}

\newpage

\subsection{نمط \textenglish{diagonal} (القطري)}

\begin{english}
\begin{lstlisting}[style=fariscode]
\FarisSetStyle{diagonal}
\FarisSetColor{RoyalBlue}
\FarisTitlePage[A][70][-][RoyalBlue][Amiri][Amiri][0.7]{@\AR{نمط قطري}@}{@\AR{تصميم بخطوط مائلة}@}
\end{lstlisting}
\end{english}

\vspace{1cm}

\noindent\textbf{\AR{النتيجة}:}\\[0.3cm]
\adjustbox{scale=0.3,center}{%
\begin{minipage}{\textwidth}
\FarisSetStyle{diagonal}
\FarisSetColor{RoyalBlue}
\FarisTitlePage[A][70][-][RoyalBlue][Amiri][Amiri][0.7]{نمط قطري}{تصميم بخطوط مائلة}
\end{minipage}
}

\vspace{6cm}

\subsection{نمط \textenglish{wave} (الموجي)}

\begin{english}
\begin{lstlisting}[style=fariscode]
\FarisSetStyle{wave}
\FarisSetColor{Teal}
\FarisTitlePage[A][75][-][Teal][Amiri][Amiri][0.65]{@\AR{نمط موجي}@}{@\AR{تصميم بموجات منحنية}@}
\end{lstlisting}
\end{english}

\vspace{7cm}

\noindent\textbf{\AR{النتيجة}:}\\[0.3cm]
\adjustbox{scale=0.3,center}{%
\begin{minipage}{\textwidth}
\FarisSetStyle{wave}
\FarisSetColor{Teal}
\FarisTitlePage[A][75][-][Teal][Amiri][Amiri][0.65]{نمط موجي}{تصميم بموجات منحنية}
\end{minipage}
}

\newpage

\subsection{نمط \textenglish{badge} (الشارة)}

\begin{english}
\begin{lstlisting}[style=fariscode]
\FarisSetStyle{badge}
\FarisSetColor{Purple}
\FarisTitlePage[A][65][-][Purple][Amiri][Amiri][0.8]{@\AR{نمط الشارة}@}{@\AR{تصميم دائري متداخل}@}
\end{lstlisting}
\end{english}

\vspace{1cm}

\noindent\textbf{\AR{النتيجة}:}\\[0.3cm]
\adjustbox{scale=0.3,center}{%
\begin{minipage}{\textwidth}
\FarisSetStyle{badge}
\FarisSetColor{Purple}
\FarisTitlePage[A][65][-][Purple][Amiri][Amiri][0.8]{نمط الشارة}{تصميم دائري متداخل}
\end{minipage}
}

\vspace{6cm}

\subsection{نمط \textenglish{diamond} (الماسي البسيط)}

\begin{english}
\begin{lstlisting}[style=fariscode]
\FarisSetStyle{diamond}
\FarisSetColor{OrangeRed}
\FarisTitlePage[A][60][-][OrangeRed][Amiri][Amiri][0.75]{@\AR{نمط ماسي}@}{@\AR{تصميم مربع مدور}@}
\end{lstlisting}
\end{english}

\vspace{6cm}

\noindent\textbf{\AR{النتيجة}:}\\[0.3cm]
\adjustbox{scale=0.3,center}{%
\begin{minipage}{\textwidth}
\FarisSetStyle{diamond}
\FarisSetColor{OrangeRed}
\FarisTitlePage[A][60][-][OrangeRed][Amiri][Amiri][0.75]{نمط ماسي}{تصميم مربع مدور}
\end{minipage}
}

\newpage

\subsection{نمط \textenglish{minimal} (البسيط)}

\begin{english}
\begin{lstlisting}[style=fariscode]
\FarisSetStyle{minimal}
\FarisSetColor{DarkSlateGray}
\FarisTitlePage[A][55][-][DarkSlateGray][Amiri][Amiri][0.9]{@\AR{نمط بسيط}@}{@\AR{تصميم بحد أدنى من العناصر}@}
\end{lstlisting}
\end{english}

\vspace{1cm}

\noindent\textbf{\AR{النتيجة}:}\\[0.3cm]
\adjustbox{scale=0.3,center}{%
\begin{minipage}{\textwidth}
\FarisSetStyle{minimal}
\FarisSetColor{DarkSlateGray}
\FarisTitlePage[A][55][-][DarkSlateGray][Amiri][Amiri][0.9]{نمط بسيط}{تصميم بحد أدنى من العناصر}
\end{minipage}
}

\vspace{6cm}

\subsection{نمط \textenglish{ribbon} (الشريط)}

\begin{english}
\begin{lstlisting}[style=fariscode]
\FarisSetStyle{ribbon}
\FarisSetColor{Maroon}
\FarisTitlePage[A][70][-][Maroon][Amiri][Amiri][0.6]{@\AR{نمط الشريط}@}{@\AR{تصميم بأشرطة متقاطعة}@}
\end{lstlisting}
\end{english}

\vspace{6cm}

\noindent\textbf{\AR{النتيجة}:}\\[0.3cm]
\adjustbox{scale=0.3,center}{%
\begin{minipage}{\textwidth}
\FarisSetStyle{ribbon}
\FarisSetColor{Maroon}
\FarisTitlePage[A][70][-][Maroon][Amiri][Amiri][0.6]{نمط الشريط}{تصميم بأشرطة متقاطعة}
\end{minipage}
}

\newpage

\subsection{نمط \textenglish{sidebar} (الشريط الجانبي)}

\begin{english}
\begin{lstlisting}[style=fariscode]
\FarisSetStyle{sidebar}
\FarisSetColor{Navy}
\FarisTitlePage[A][65][-][Navy][Amiri][Amiri][0.7]{@\AR{نمط الشريط الجانبي}@}{@\AR{تصميم بشريط يساري}@}
\end{lstlisting}
\end{english}

\vspace{1cm}

\noindent\textbf{\AR{النتيجة}:}\\[0.3cm]
\adjustbox{scale=0.3,center}{%
\begin{minipage}{\textwidth}
\FarisSetStyle{sidebar}
\FarisSetColor{Navy}
\FarisTitlePage[A][65][-][Navy][Amiri][Amiri][0.7]{نمط الشريط الجانبي}{تصميم بشريط يساري}
\end{minipage}
}

\vspace{6cm}

\subsection{نمط \textenglish{diamondx} المتقدم — أساسي}

\begin{english}
\begin{lstlisting}[style=fariscode]
\FarisSetStyle{diamondx}
\FarisSetColor{DeepSkyBlue}
\FarisTitlePage[A][75][-][DeepSkyBlue][Amiri][Amiri][0.6]{@\AR{النمط الماسي المتقدم}@}{@\AR{إصدار محسّن من النمط الماسي}@}
\end{lstlisting}
\end{english}

\vspace{6cm}

\noindent\textbf{\AR{النتيجة}:}\\[0.3cm]
\adjustbox{scale=0.3,center}{%
\begin{minipage}{\textwidth}
\FarisSetStyle{diamondx}
\FarisSetColor{DeepSkyBlue}
\FarisTitlePage[A][75][-][DeepSkyBlue][Amiri][Amiri][0.6]{النمط الماسي المتقدم}{إصدار محسّن من النمط الماسي}
\end{minipage}
}

\newpage

\subsection{نمط \textenglish{diamondx} مع الألوان المتممة}

\begin{english}
\begin{lstlisting}[style=fariscode]
\FarisSetStyle{diamondx}
\FarisSetColor{Red}
\FarisDiamondBackComplementOn    % @\AR{تفعيل الألوان المتممة}@
\FarisTitlePage[A][70][-][Red][Amiri][Amiri][0.65]{@\AR{النمط مع ألوان متممة}@}{@\AR{استخدام الألوان المعاكسة}@}
\end{lstlisting}
\end{english}

\vspace{1cm}

\noindent\textbf{\AR{النتيجة}:}\\[0.3cm]
\adjustbox{scale=0.3,center}{%
\begin{minipage}{\textwidth}
\FarisSetStyle{diamondx}
\FarisSetColor{Red}
\FarisDiamondBackComplementOn
\FarisTitlePage[A][70][-][Red][Amiri][Amiri][0.65]{النمط مع ألوان متممة}{استخدام الألوان المعاكسة}
\end{minipage}
}

\vspace{6cm}

\subsubsection{نمط \textenglish{diamondx} مع تخصيص التينتات}

\begin{english}
\begin{lstlisting}[style=fariscode]
\FarisSetStyle{diamondx}
\FarisSetColor{Green}
\FarisDiamondBackComplementOff   % @\AR{إلغاء الألوان المتممة}@
\FarisSetDiamondTints{40}{25}    % @\AR{تينت الأمامي 40 والخلفي A بـ25}@
\FarisSetDiamondBackTintII{35}   % @\AR{تينت الخلفي B بـ35}@
\FarisTitlePage[A][68][-][Green][Amiri][Amiri][0.7]{@\AR{نمط بتينتات مخصصة}@}{@\AR{تحكم في شدة الألوان}@}
\end{lstlisting}
\end{english}

\vspace{7cm}

\noindent\textbf{\AR{النتيجة}:}\\[0.3cm]
\adjustbox{scale=0.3,center}{%
\begin{minipage}{\textwidth}
\FarisSetStyle{diamondx}
\FarisSetColor{Green}
\FarisDiamondBackComplementOff
\FarisSetDiamondTints{40}{25}
\FarisSetDiamondBackTintII{35}
\FarisTitlePage[A][68][-][Green][Amiri][Amiri][0.7]{نمط بتينتات مخصصة}{تحكم في شدة الألوان}
\end{minipage}
}

\newpage

\subsubsection{نمط \textenglish{diamondx} مع تخصيص الزوايا}

\begin{english}
\begin{lstlisting}[style=fariscode]
\FarisSetStyle{diamondx}
\FarisSetColor{Violet}
\FarisSetDiamondAnglePair{30}    % @\AR{زوايا ±30°}@
\FarisSetDiamondBackDelta{20}    % @\AR{فرق 20° بين الخلفيين}@
\FarisTitlePage[A][72][-][Violet][Amiri][Amiri][0.6]{@\AR{نمط بزوايا مخصصة}@}{@\AR{تدوير المربعات بزوايا مختلفة}@}
\end{lstlisting}
\end{english}

\vspace{1cm}

\noindent\textbf{\AR{النتيجة}:}\\[0.3cm]
\adjustbox{scale=0.3,center}{%
\begin{minipage}{\textwidth}
\FarisSetStyle{diamondx}
\FarisSetColor{Violet}
\FarisSetDiamondAnglePair{30}
\FarisSetDiamondBackDelta{20}
\FarisTitlePage[A][72][-][Violet][Amiri][Amiri][0.6]{نمط بزوايا مخصصة}{تدوير المربعات بزوايا مختلفة}
\end{minipage}
}

\vspace{6cm}

\subsubsection{نمط \textenglish{diamondx} مع الألوان الثلاثية}

\begin{english}
\begin{lstlisting}[style=fariscode]
\FarisSetStyle{diamondx}
\FarisSetColor{Orange}
\FarisSetDiamondTriadOffset{120} % @\AR{إزاحة ثلاثية 120°}@
\FarisSetDiamondTints{30}{30}
\FarisSetDiamondBackTintII{40}
\FarisTitlePage[A][66][-][Orange][Amiri][Amiri][0.75]{@\AR{نمط بألوان ثلاثية}@}{@\AR{استخدام عجلة الألوان}@}
\end{lstlisting}
\end{english}

\vspace{6.5cm}

\noindent\textbf{\AR{النتيجة}:}\\[0.3cm]
\adjustbox{scale=0.3,center}{%
\begin{minipage}{\textwidth}
\FarisSetStyle{diamondx}
\FarisSetColor{Orange}
\FarisSetDiamondTriadOffset{120}
\FarisSetDiamondTints{30}{30}
\FarisSetDiamondBackTintII{40}
\FarisTitlePage[A][66][-][Orange][Amiri][Amiri][0.75]{نمط بألوان ثلاثية}{استخدام عجلة الألوان}
\end{minipage}
}

\newpage

\subsubsection{تخصيص كامل للنمط \textenglish{diamondx}}

\begin{english}
\begin{lstlisting}[style=fariscode]
\FarisSetStyle{diamondx}
\FarisSetColor{DarkMagenta}
\FarisDiamondBackComplementOn
\FarisSetDiamondTints{25}{35}
\FarisSetDiamondBackTintII{45}
\FarisSetDiamondAnglePair{60}
\FarisSetDiamondBackDelta{25}
\FarisSetDiamondTriadOffset{90}
\FarisTitlePage[A][64][-][DarkMagenta][Amiri][Amiri][0.8]{@\AR{تخصيص شامل}@}{@\AR{جميع الإعدادات مجتمعة}@}
\end{lstlisting}
\end{english}

\vspace{1.5cm}

\noindent\textbf{\AR{النتيجة}:}\\[0.3cm]
\adjustbox{scale=0.3,center}{%
\begin{minipage}{\textwidth}
\FarisSetStyle{diamondx}
\FarisSetColor{DarkMagenta}
\FarisDiamondBackComplementOn
\FarisSetDiamondTints{25}{35}
\FarisSetDiamondBackTintII{45}
\FarisSetDiamondAnglePair{60}
\FarisSetDiamondBackDelta{25}
\FarisSetDiamondTriadOffset{90}
\FarisTitlePage[A][64][-][DarkMagenta][Amiri][Amiri][0.8]{تخصيص شامل}{جميع الإعدادات مجتمعة}
\end{minipage}
}

\vspace{5cm}

% ==========================  section = eight  =========================
\section{استخدام خطوط إنجليزية}
% ==================================================================

\begin{english}
\begin{lstlisting}[style=fariscode]
\FarisSetStyle{badge}
\FarisSetColor{SteelBlue}
\FarisTitlePage[E][80][-][SteelBlue][Times New Roman][Georgia][0.7]{Advanced LaTeX Package}{For Beautiful Title Pages}
\end{lstlisting}
\end{english}

\vspace{6.5cm}

\noindent\textbf{\AR{النتيجة}:}\\[0.3cm]
\adjustbox{scale=0.3,center}{%
\begin{minipage}{\textwidth}
\FarisSetStyle{badge}
\FarisSetColor{SteelBlue}
\FarisTitlePage[E][70][-][SteelBlue][Times New Roman][Georgia][0.7]{Advanced LaTeX Package}{For Beautiful Title Pages}
\end{minipage}
}

\newpage

% ==========================  section = nine  =========================
\section{مثال مختلط: عربي وإنجليزي}
% ==================================================================

\begin{english}
\begin{lstlisting}[style=fariscode]
\FarisSetStyle{wave}
\FarisSetColor{Chocolate}
\FarisTitlePage[A][85][-][Chocolate][Amiri][Amiri][0.5]{@\AR{دليل المستخدم}@}{User Manual - @\AR{الإصدار الثاني}@}
\end{lstlisting}
\end{english}

\vspace{1cm}

\noindent\textbf{\AR{النتيجة}:}\\[0.3cm]
\adjustbox{scale=0.3,center}{%
\begin{minipage}{\textwidth}
\FarisSetStyle{wave}
\FarisSetColor{Chocolate}
\FarisTitlePage[A][85][-][Chocolate][Amiri][Amiri][0.5]{دليل المستخدم}{User Manual - الإصدار الثاني}
\end{minipage}
}

\vspace{6cm}

% ==========================  section = ten  ===========================
\section{استخدام لون مخصص بصيغة \textenglish{RGB}}
% ==================================================================

\begin{english}
\begin{lstlisting}[style=fariscode]
\definecolor{CustomBlue}{RGB}{25,118,210}
\FarisSetStyle{diagonal}
\FarisSetColor{CustomBlue}
\FarisTitlePage[A][78][-][CustomBlue][Amiri][Amiri][0.65]{@\AR{لون مخصص}@}{@\AR{استخدام صيغة RGB}@}
\end{lstlisting}
\end{english}

\vspace{6cm}

\noindent\textbf{\AR{النتيجة}:}\\[0.3cm]
\adjustbox{scale=0.3,center}{%
\begin{minipage}{\textwidth}
\definecolor{CustomBlue}{RGB}{25,118,210}
\FarisSetStyle{diagonal}
\FarisSetColor{CustomBlue}
\FarisTitlePage[A][78][-][CustomBlue][Amiri][Amiri][0.65]{لون مخصص}{استخدام صيغة RGB}
\end{minipage}
}

\newpage

% ==========================  section = eleven ========================
\section{ملخص الأوامر والمعاملات}
% ==================================================================

\subsection{بناء الأمر الأساسي}

\begin{english}
\begin{lstlisting}[style=fariscode]
\FarisTitlePage[Lang][TitlePt][SubtitlePt][Color][MainFont][SubFont][SubRatio]{Title}{Subtitle}
\end{lstlisting}
\end{english}

\subsection{المعاملات الاختيارية}

\begin{itemize}
  \item \textbf{Lang}: \AR{اللغة —} A \AR{للعربية،} E \AR{للإنجليزية}
  \item \textbf{TitlePt}: \AR{حجم العنوان الرئيسي بالنقاط (0 أو - = افتراضي)}
  \item \textbf{SubtitlePt}: \AR{حجم العنوان الفرعي بالنقاط (0 أو - = تلقائي)}
  \item \textbf{Color}: \AR{لون التصميم العام}
  \item \textbf{MainFont}: \AR{خط العنوان الرئيسي}
  \item \textbf{SubFont}: \AR{خط العنوان الفرعي}
  \item \textbf{SubRatio}: \AR{ نسبة حجم الفرعي من الرئيسي \en{ (1.0 — 0.1)}}
\end{itemize}

\subsection{أوامر التخصيص العامة}

\begin{english}
\begin{lstlisting}[style=fariscode]
\FarisSetStyle{style}                    % @\AR{تحديد النمط}@
\FarisSetColor{color}                    % @\AR{تحديد اللون الأساسي}@
\FarisSetAccentColor{color}              % @\AR{تحديد لون الإبراز}@
\FarisSetDims{width}{height}{corner}     % @\AR{تخصيص الأبعاد}@
\end{lstlisting}
\end{english}

\subsection{أوامر تلوين النصوص — \textenglish{v1.1}}

\begin{english}
\begin{lstlisting}[style=fariscode]
\FarisSetTitleTextColor{color}     % @\AR{لون العنوان الرئيسي}@
\FarisSetSubtitleTextColor{color}  % @\AR{لون العنوان الفرعي}@
\FarisResetTitleTextColor          % @\AR{إعادة لون الرئيسي للافتراضي}@
\FarisResetSubtitleTextColor       % @\AR{إعادة لون الفرعي للافتراضي}@
\FarisResetTextColors              % @\AR{إعادة كليهما للافتراضي}@
\end{lstlisting}
\end{english}

\subsection{أوامر خاصة بالنمط \textenglish{diamondx}}

\begin{english}
\begin{lstlisting}[style=fariscode]
\FarisDiamondBackComplementOn              % @\AR{تفعيل الألوان المتممة}@
\FarisDiamondBackComplementOff             % @\AR{إلغاء الألوان المتممة}@
\FarisSetDiamondTints{front}{backA}        % @\AR{تينت الأمامي والخلفي A}@
\FarisSetDiamondBackTintII{backB}          % @\AR{تينت الخلفي B}@
\FarisSetDiamondAnglePair{angle}           % @\AR{زوايا متقابلة ±angle}@
\FarisSetDiamondAngles{backAngle}{frontAngle} % @\AR{زوايا منفصلة}@
\FarisSetDiamondBackDelta{delta}           % @\AR{فرق الزاوية بين الخلفيين}@
\FarisSetDiamondTriadOffset{offset}        % @\AR{إزاحة الألوان الثلاثية}@
\end{lstlisting}
\end{english}

\newpage
 
% ==========================  section = twelve =========================
\section{تعدد الأسطر في العناوين}
% ==================================================================
 
\AR{تدعم الحزمة الكسر التلقائي والكسر اليدوي للعناوين الطويلة، مع الحفاظ
على اتجاه النص العربي من اليمين إلى اليسار في جميع الأسطر.}
 
\subsection{الكسر اليدوي بـ \textenglish{\textbackslash\textbackslash}}
 
\begin{english}
\begin{lstlisting}[style=fariscode]
% @\AR{كسر السطر يدوياً في العنوان الفرعي}@
\FarisTitlePage[A][90][-][Chocolate][Diwani Simple Striped][Amiri][0.7]%
  {@\AR{الفصل السابع}@}{@\AR{الكيمياء العضوية \textbackslash\textbackslash \, الهيدروكاربونات المعوضة}@}
\end{lstlisting}
\end{english}
 
\vspace{1.5cm}
 
\noindent\textbf{\AR{النتيجة}:}\\[0.3cm]
\adjustbox{scale=0.3,center}{%
\begin{minipage}{\textwidth}
\FarisResetTextColors
\FarisTitlePage[A][90][-][Chocolate][Diwani Simple Striped][Amiri][0.7]%
  {الفصل السابع}{الكيمياء العضوية \\ الهيدروكاربونات المعوضة}
\end{minipage}
}
 
\vspace{5cm}
 
\subsection{الكسر التلقائي للنص الطويل}
 
\begin{english}
\begin{lstlisting}[style=fariscode]
% @\AR{النص الطويل يُكسر تلقائياً عند امتلاء السطر (82\% من عرض الصفحة)}@
\FarisTitlePage[A][90][-][NavyBlue][Amiri][Amiri][0.6]%
  {@\AR{الفصل الثامن}@}{@\AR{الميكانيكا الكلاسيكية وقوانين نيوتن في الحركة والقوى}@}
\end{lstlisting}
\end{english}
 
\vspace{6.5cm}
 
\noindent\textbf{\AR{النتيجة}:}\\[0.3cm]
\adjustbox{scale=0.3,center}{%
\begin{minipage}{\textwidth}
\FarisResetTextColors
\FarisTitlePage[A][90][-][NavyBlue][Amiri][Amiri][0.6]%
  {الفصل الثامن}{الميكانيكا الكلاسيكية وقوانين نيوتن في الحركة والقوى}
\end{minipage}
}
 
\vspace{4cm}
 
\subsection{كسر السطر في العنوان الرئيسي والفرعي معاً}
 
\begin{english}
\begin{lstlisting}[style=fariscode]
\FarisTitlePage[A][90][-][DarkOrange1][Diwani Simple Striped][Amiri][0.6]%
  {@\AR{الفصل التاسع \textbackslash\textbackslash \, مقدمة في الحركيات}@}{@\AR{الكيمياء الفيزيائية \textbackslash\textbackslash \, سرعة التفاعلات}@}
\end{lstlisting}
\end{english}
 
\vspace{1cm}
 
\noindent\textbf{\AR{النتيجة}:}\\[0.3cm]
\adjustbox{scale=0.3,center}{%
\begin{minipage}{\textwidth}
\FarisResetTextColors
\FarisTitlePage[A][90][-][DarkOrange1][Diwani Simple Striped][Amiri][0.6]%
  {الفصل التاسع \\ مقدمة في الحركيات}{الكيمياء الفيزيائية \\ سرعة التفاعلات}
\end{minipage}
}

\vspace{7cm}

% ==================================================================
\section{الأنماط المتاحة}
% ==================================================================

\begin{itemize}
  \item \textbf{triple}: \AR{ثلاثي الطبقات (افتراضي)}
  \item \textbf{diagonal}: \AR{قطري}
  \item \textbf{wave}: \AR{موجي}
  \item \textbf{badge}: \AR{شارة دائرية}
  \item \textbf{diamond}: \AR{ماسي بسيط}
  \item \textbf{minimal}: \AR{بسيط}
  \item \textbf{ribbon}: \AR{شريطي}
  \item \textbf{sidebar}: \AR{شريط جانبي}
  \item \textbf{diamondx}: \AR{ماسي متقدم}
\end{itemize}

\newpage

% ==================================================================
\section{نصائح الاستخدام}
% ==================================================================

\begin{itemize}
  \item \AR{استخدم XeLaTeX أو LuaLaTeX للتجميع}
  \item \AR{تأكد من تثبيت الخطوط المطلوبة}
  \item \AR{استخدم الترميز UTF-8}
  \item \AR{اكتب \texttt{\textbackslash FarisResetTextColors} قبل كل غلاف جديد إذا غيّرت الألوان}
  \item \AR{القيمة 0 في TitlePt أو SubtitlePt تُعامَل كـ ``-'' وتُشغِّل الحساب التلقائي}
  \item \AR{الفرعي لا يتجاوز الرئيسي أبداً — الحارس التلقائي يُطبِّق 4/5 إذا تجاوز}
  \item \AR{جرّب الأنماط المختلفة لاختيار الأنسب لمشروعك}
  \item \AR{اضبط الألوان والخطوط حسب طبيعة المادة الدراسية}
\end{itemize}


% ==================================================================
\section{  كلمة المؤلف}
% ==================================================================



  هذه الحزمة وُلدت من حاجة حقيقية: كنت أبحث عن طريقة بسيطة
  وأنيقة لتصميم صفحات العناوين في مستنداتي العربية باستخدام
\en{ XeLaTeX} ، فلم أجد ما يُرضيني، فقررت أن أبنيها بنفسي.

  بدأت بفكرة واحدة، ثم تطورت عبر جلسات طويلة من التجريب والتصحيح والتفكير، حتى أصبحت ما تراه اليوم: تسعة أنماط، نظام ألوان مرن، حساب تلقائي للأحجام، وحراس أمان تمنع أي خطأ من الوصول إلى الناتج النهائي.

  كل سطر في هذه الحزمة كُتب بعناية، وكل حارس أُضيف بعد تجربة فاشلة تعلمت منها. إن وجدت فيها فائدة، فهذا يكفيني. وإن وجدت فيها خطأ، فأنا أول من يريد أن يعرف.

  أتمنى أن تُزيّن هذه الحزمة مستنداتك وكتبك ورسائلك، وأن تجعل البداية دائماً جميلة.
\\
\begin{center}
                                  فارس  الحيدري \\
                                 \en{  classsec81@gmail.com} \\
                                  \en{2026/03/21 — v1.1} 
\end{center}

\end{document}
